% ECONOMICS_PROBLEM: {"id": "D74-356", "title": "Conflict, fragility and poverty traps", "jel_code": "D74", "source_id": "OP-0039"}
% !TeX program = xelatex
\documentclass[11pt,a4paper]{article}
\usepackage[margin=23mm]{geometry}
\usepackage{fontspec}
\setmainfont{Latin Modern Roman}
\usepackage{amsmath,amssymb,mathtools}
\usepackage{xcolor,enumitem,booktabs,longtable,array,xurl,hyperref}
\definecolor{muted}{HTML}{53636C}
\hypersetup{colorlinks=true,urlcolor=blue,linkcolor=blue}
\setlength{\parindent}{0pt}
\setlength{\parskip}{5pt}
\newcommand{\R}{\mathbb R}
\newcommand{\E}{\mathbb E}
\newcommand{\N}{\mathbb N}
\newcommand{\ind}{\mathbf 1}
\newcommand{\Var}{\operatorname{Var}}
\newcommand{\Cov}{\operatorname{Cov}}
\newcommand{\supp}{\operatorname{supp}}
\DeclareMathOperator*{\argmin}{arg\,min}
\DeclareMathOperator*{\argmax}{arg\,max}
\newcommand{\entrymeta}[3]{{\sffamily\small\color{muted}\textbf{\texttt{#1}}\quad|\quad #2\par #3\par}}
\newcommand{\Problemref}[1]{\texttt{#1}}
\newcommand{\JELref}[2]{\texttt{#2} (JEL \texttt{#1})}
\begin{document}
\section*{Conflict, fragility and poverty traps}
\textbf{Problem ID:} \texttt{D74-356}\quad \textbf{JEL:} \texttt{D74}\par
% BEGIN ECONOMICS STATEMENT
\label{problem:OP-0039}
\label{atlas:prob:D9}
\noindent\textbf{Original atlas ID:} D9 (entry 39). \textbf{Classification:} Editorial reconstruction sequencing and stochastic durability problem.

\noindent\textbf{Original atlas question:} Which interventions move post-conflict economies to a durable higher-capacity equilibrium?

\subsection*{Definitions and assumptions}
Model a post-conflict economy by state $S_t=(K_t,C_t,V_t,Z_t)$: productive capital, administrative capability, violence, and political/network conditions. A reconstruction rule $u$ selects security, transfers, and services using observed histories, subject to budget and political feasibility. Model $M$ specifies controlled transition probabilities, actors' incentives, and an equilibrium selection rule. Define a measurable target region $G$ by announced minimum income/capability levels and a maximum violence threshold. Let $\tau_G=\inf\{t\geq0:S_t\in G\}$, with $\inf\varnothing=\infty$. Do not assume that a poverty trap exists merely because poverty and conflict are correlated.
\subsection*{Formal research question}
For an entry horizon $H$ and durability horizon $J$, identify or bound
\[
 p_M(u)=\Pr_M^u\!\left(\tau_G\leq H,
          \ S_t\in G\text{ for every integer }t\in[\tau_G,\tau_G+J]\right).
\]
Compare feasible rules using both $p_M(u)-p_M(u_0)$ and discounted welfare net of real reconstruction costs. Determine whether alternative sequences of security, transfers, and services have distinguishable causal effects, including spillovers to neighboring regions. If a stationary long-run model is used, establish existence of its invariant distribution and characterize the policy effect on mass in $G$; multiple equilibria or trapping classes must be demonstrated within that model.
\subsection*{Interpretation and scope}
The displayed probability is a finite-horizon durability target, not a claim of permanent peace. In a stochastic economy, a deterministic stable steady state, an absorbing class, and a highly persistent low-income region are different concepts. Migration can improve measured local income by changing who remains; baseline-cohort outcomes should accompany regional measures. Reconstruction assignment is often endogenous to security and political control. Historical or weather instruments need their own exclusion restrictions and cannot automatically identify modern reconstruction packages. Randomized sequencing identifies only implemented sequences with adequate support and a specified spillover design. Noncompliance, displacement, and conflict reporting can themselves respond to policy and require explicit measurement assumptions.
\subsection*{Source audit and status}
[S1] studies economic shocks and conflict, [S2] reviews civil war, and [S3] studies historical state capacity and collective action in Vietnam. [S3] is not a randomized post-conflict reconstruction evaluation; its relevance is indirect. The sequencing-and-durability problem is an editorial extension. These sources neither prove a universal poverty trap nor certify a generally unresolved 2026 intervention problem.

\subsection*{References}
\begin{enumerate}[label=\textnormal{[S\arabic*]},leftmargin=*]
\item Edward Miguel, Shanker Satyanath, and Ernest Sergenti (2004). \emph{Economic Shocks and Civil Conflict: An Instrumental Variables Approach}. Journal of Political Economy 112(4), 725--753. \url{https://doi.org/10.1086/421174}. \textit{Locator:} Primary publisher, working-paper, or author record: bibliographic information and abstract; full-text theorem verification is not claimed. \textit{Role:} Rainfall/economic-shock evidence on conflict; not identification of reconstruction-policy sequencing.
\item Christopher Blattman and Edward Miguel (2010). \emph{Civil War}. Journal of Economic Literature 48(1), 3--57. \url{https://doi.org/10.1257/jel.48.1.3}. \textit{Locator:} Primary publisher, working-paper, or author record: bibliographic information and abstract; full-text theorem verification is not claimed. \textit{Role:} Survey of conflict theory and evidence; background for dynamic political models.
\item Melissa Dell, Nathan Lane, and Pablo Querubin (2018). \emph{The Historical State, Local Collective Action, and Economic Development in Vietnam}. Econometrica 86(6), 2083--2121. \url{https://doi.org/10.3982/ECTA15122}. \textit{Locator:} Primary publisher, working-paper, or author record: bibliographic information and abstract; full-text theorem verification is not claimed. \textit{Role:} Historical state-capacity and collective-action evidence; indirect motivation for post-conflict capacity, not a reconstruction trial.
\end{enumerate}

\subsection*{Common conventions: Source mathematical conventions}

Unless an entry states otherwise, agent and firm sets are finite. State and action spaces are measurable, strategies and outcome maps are measurable, probability laws are specified, and expectations exist for the targets under discussion. Infinite-horizon discount factors lie in $(0,1)$ and discounted objectives are finite. Norms, tolerances and complexity input sizes must be specified for computational comparisons. Constants or primitives that appear locally are inputs, not empirical facts asserted by this catalogue.

\textbf{Identification.} A statistical model $m\in\mathcal M$ implies an observable law $P_m$ and target $\theta(m)$. At law $P$, its identified set is
\[
\Theta(P;\mathcal M)=\{\theta(m):m\in\mathcal M,\ P_m=P\}.
\]
Point identification holds when this set is a singleton, or equivalently when $P_{m_1}=P_{m_2}$ implies $\theta(m_1)=\theta(m_2)$ for all compatible models. Sharp bounds mean bounds attained, or approached when the boundary is not attained, by compatible models. Writing this set is a definition, not a solution: the substantive task is to establish informative restrictions and derive or estimate the set. An unrestricted-confounding model may give only trivial bounds.

\textbf{Inference.} Identification, estimability and valid uncertainty quantification are separate questions. Fix $\alpha\in(0,1)$. A confidence procedure $C_n$ over a declared law class $\mathcal P$ has asymptotic uniform coverage at level $1-\alpha$ when
\[
\liminf_{n\to\infty}\inf_{P\in\mathcal P}\Pr_P\{\theta(P)\in C_n\}\geq1-\alpha.
\]
For set-identified targets the coverage event must be defined separately, for example coverage of the full identified set. If an entry uses identification-indexed models instead of uniquely indexed laws, the same coverage convention is applied over those models. Pointwise limits do not establish uniform validity, and asymptotic validity does not guarantee finite-sample precision. Sample sizes, dependence structures and weak-identification sequences are part of each application.

\textbf{Counterfactuals.} $Y_i(a)$ is a potential outcome under a specified intervention, including its timing, implementation and versions. It is not $Y_i$ conditional on voluntarily choosing $a$. A full intervention vector permits interference; shorthand scalar treatment notation does not silently impose no interference. Assignment, selection, attrition, anticipation and measurement must be specified. A claim about an instrument's local effect is not automatically a population policy effect.

\textbf{Economic equilibrium.} A counterfactual model includes best responses, constraints, feasibility, market clearing, state transitions and expectations consistent with the stated information. When several equilibria exist, either a declared selector is an input or the target is set-valued. Comparisons across policy regimes must use the stated selection convention. Reduced-form relationships are not presumed invariant after an equilibrium-changing intervention.

\textbf{Optimization and welfare.} Optimization is over the stated feasible set. If attainment is not established, $\sup$ or $\inf$ and $\varepsilon$-optimal choices replace an asserted maximizer or minimizer. For example, with value $V=\sup_{a\in\mathcal A}W(a)<\infty$, an $\varepsilon$-optimal set is $\{a\in\mathcal A:W(a)\geq V-\varepsilon\}$ for $\varepsilon>0$. Benefits, costs, risk and health or environmental outcomes can be aggregated only using declared units and conversion weights. Interpersonal utility weights, the treatment of fiscal transfers and foreign profits, discounting, and the behavioral welfare criterion are explicit inputs. A comparative-static derivative additionally requires existence and appropriate differentiability of the selected counterfactual.

\textbf{Sources.} Local labels [S1], [S2], and so forth restart in every question. Locators identify the relevant abstract, model, section or result at the level actually checked; a bibliographic or abstract-level verification is not represented as inspection of an entire proof. Working papers are distinguished from later publications and changing versions. Source summaries are paraphrased. All URL checks and editorial formulations refer to the audit date stated above.
% END ECONOMICS STATEMENT
\end{document}
